\documentclass{article}
\usepackage{amsmath,amssymb,amsthm}
\usepackage{algorithm}
\usepackage{algpseudocode}
\usepackage{bm}
\usepackage{graphicx}
\usepackage{hyperref}
\usepackage{enumitem}
\newtheorem{theorem}{Theorem}
\newtheorem{lemma}{Lemma}
\newtheorem{assumption}{Assumption}
\newtheorem{corollary}{Corollary}
\newcommand{\E}{\mathbb{E}}
\newcommand{\R}{\mathbb{R}}
\newcommand{\norm}[1]{\left\lVert#1\right\rVert}
\title{Stochastic Gradient Langevin Dynamics \\ with Decaying Noise \\ (SGLD‑D)}
\author{}
\date{}
\begin{document}
\maketitle

\section*{Algorithm Name}
\textbf{Stochastic Gradient Langevin Dynamics with Decaying Noise (SGLD‑D)}  

The method combines a standard stochastic gradient step with an additive Gaussian perturbation whose variance is scheduled to decay over iterations. It is a discretisation of the overdamped Langevin diffusion for sampling from a (log‑concave) target distribution, but here the decay of the temperature is exploited to obtain optimisation guarantees for convex smooth objectives.

\section*{Mathematical Formulation}
Let $f:\R^{d}\to\R$ be the objective.  At iteration $k\ge 0$ the algorithm produces a random iterate $x_{k}$ according to  

\[
\boxed{
x_{k+1}=x_{k}-\alpha_{k}\,g_{k}+\sqrt{2\alpha_{k}\,\tau_{k}}\;\xi_{k}}
\tag{1}
\]

where  

\begin{itemize}[nosep]
\item $\alpha_{k}>0$ is the step size (learning rate);
\item $g_{k}$ is a stochastic gradient satisfying $\E[g_{k}\mid x_{k}]=\nabla f(x_{k})$;
\item $\tau_{k}\ge0$ is the temperature (noise variance) at iteration $k$;
\item $\xi_{k}\sim\mathcal N(0,\mathbf I_{d})$ are i.i.d. standard Gaussian vectors, independent of the history $\mathcal F_{k}:=\sigma(x_{0},\xi_{0},\dots,\xi_{k-1})$.
\end{itemize}

The noise term can be written equivalently as $\eta_{k}$ with $\E[\eta_{k}\mid\mathcal F_{k}]=0$ and $\E[\|\eta_{k}\|^{2}\mid\mathcal F_{k}]=2\alpha_{k}\tau_{k}d$.

\section*{Pseudocode}
\begin{algorithm}[H]
\caption{SGLD‑D}
\begin{algorithmic}[1]
\State \textbf{Input:} initial point $x_{0}\in\R^{d}$, step–size schedule $\{\alpha_{k}\}_{k\ge0}$, temperature schedule $\{\tau_{k}\}_{k\ge0}$, max iterations $K$.
\For{$k=0$ \textbf{to} $K-1$}
    \State Sample a mini‑batch (or a random datum) and compute a stochastic gradient $g_{k}$.
    \State Sample $\xi_{k}\sim\mathcal N(0,\mathbf I_{d})$.
    \State $x_{k+1}\gets x_{k}-\alpha_{k}g_{k}+\sqrt{2\alpha_{k}\tau_{k}}\;\xi_{k}$.
\EndFor
\State \textbf{Output:} $x_{K}$ (or the averaged iterate $\bar x_{K}=\frac{1}{K}\sum_{k=0}^{K-1}x_{k}$).
\end{algorithmic}
\end{algorithm}

\section*{Dry Run Example}
Consider the one‑dimensional quadratic  

\[
f(w)=(w-3)^{2},\qquad \nabla f(w)=2(w-3).
\]

We use a deterministic gradient (i.e. $g_{k}=\nabla f(x_{k})$) and set the temperature to zero for the illustration, $\tau_{k}=0$.  
Parameters: $x_{0}=0$, constant step size $\alpha_{k}=0.1$ for all $k$.

\begin{center}
\begin{tabular}{c|c|c|c}
$k$ & $x_{k}$ & $g_{k}=2(x_{k}-3)$ & $x_{k+1}=x_{k}-0.1\,g_{k}$ \\ \hline
0 & $0$ & $-6$ & $0-0.1(-6)=0.6$ \\ 
1 & $0.6$ & $2(0.6-3)=-4.8$ & $0.6-0.1(-4.8)=1.08$ \\ 
2 & $1.08$ & $2(1.08-3)=-3.84$ & $1.08-0.1(-3.84)=1.464$ \\ 
\end{tabular}
\end{center}

Corresponding objective values:

\[
f(0)=9,\quad f(0.6)=5.76,\quad f(1.08)=3.6864,\quad f(1.464)=2.3526.
\]

The iterates approach the minimiser $w^{*}=3$ while the objective decreases monotonically.

\newpage
\section*{Assumptions}
\begin{assumption}[Convexity and Smoothness]\label{ass:convex}
The function $f$ is convex and $L$‑smooth, i.e. for all $x,y\in\R^{d}$,
\[
f(y)\le f(x)+\langle\nabla f(x),y-x\rangle+\frac{L}{2}\|y-x\|^{2}.
\]
\end{assumption}

\begin{assumption}[Unbiased Stochastic Gradient]\label{ass:unbiased}
For every $k$, $\E[g_{k}\mid\mathcal F_{k}]=\nabla f(x_{k})$.
\end{assumption}

\begin{assumption}[Bounded Variance]\label{ass:variance}
There exists $\sigma^{2}>0$ such that
\[
\E\big[\|g_{k}-\nabla f(x_{k})\|^{2}\mid\mathcal F_{k}\big]\le\sigma^{2}\quad\forall k.
\]
\end{assumption}

\begin{assumption}[Step‑size and Temperature Schedules]\label{ass:schedules}
The sequences $\{\alpha_{k}\}$ and $\{\tau_{k}\}$ satisfy
\[
\alpha_{k}>0,\qquad \sum_{k=0}^{\infty}\alpha_{k}= \infty,\qquad 
\sum_{k=0}^{\infty}\alpha_{k}^{2}<\infty,
\]
\[
\tau_{k}\ge0,\qquad \tau_{k}\downarrow0,\qquad 
\sum_{k=0}^{\infty}\alpha_{k}\tau_{k}<\infty .
\]
\end{assumption}

\section*{Main Convergence Theorem}
\begin{theorem}[Expected Sub‑optimality of SGLD‑D]\label{thm:main}
Under Assumptions \ref{ass:convex}–\ref{ass:schedules}, let $x_{k}$ be generated by \eqref{1} and define the averaged iterate
\[
\bar x_{K}:=\frac{1}{K}\sum_{k=0}^{K-1}x_{k}.
\]
Then for any $K\ge1$,
\[
\E\!\big[f(\bar x_{K})-f(x^{\star})\big]
\;\le\;
\frac{\|x_{0}-x^{\star}\|^{2}}{2\sum_{k=0}^{K-1}\alpha_{k}}
\;+\;
\frac{L\sigma^{2}}{2K}\sum_{k=0}^{K-1}\alpha_{k}
\;+\;
\frac{L d}{K}\sum_{k=0}^{K-1}\alpha_{k}\tau_{k},
\tag{2}
\]
where $x^{\star}$ denotes any minimiser of $f$.
In particular, choosing $\alpha_{k}=c/\sqrt{k+1}$ and $\tau_{k}=c_{\tau}/(k+1)$ yields
\[
\E\!\big[f(\bar x_{K})-f(x^{\star})\big]=\mathcal O\!\left(\frac{1}{\sqrt{K}}\right).
\]
\end{theorem}

\section*{Theoretical Properties}
\begin{itemize}[nosep]
\item \textbf{Stability.} The decreasing temperature $\tau_{k}$ guarantees that the stochastic perturbation vanishes asymptotically, preventing divergence caused by persistent noise.
\item \textbf{Hyperparameter Sensitivity.} The convergence bound depends linearly on $L$, $\sigma^{2}$ and the dimension $d$ via the noise term. Over‑aggressive step sizes break the condition $\sum\alpha_{k}^{2}<\infty$ and may lead to divergence.
\item \textbf{Comparison to Baselines.}
    \begin{enumerate}[nosep]
        \item With $\tau_{k}\equiv0$ the method reduces to standard SGD; the bound (2) collapses to the classical $\mathcal O(1/\sqrt{K})$ rate under the same schedules.
        \item With constant $\tau_{k}>0$ the algorithm behaves like SGLD for sampling; the last term in (2) does not vanish, yielding a stationary bias proportional to $\tau$.
    \end{enumerate}
\end{itemize}

\section*{Discussion}
The theorem shows that, despite the presence of Gaussian perturbations, a suitably decaying temperature schedule preserves the optimisation‑type convergence guarantees of SGD for convex smooth objectives. The additional term $L d\sum\alpha_{k}\tau_{k}/K$ quantifies the price paid for exploration; it disappears as $\tau_{k}\to0$ fast enough (e.g. $\tau_{k}=O(1/k)$).

\newpage
\section*{Preliminaries (Definitions and Lemmas)}
\begin{lemma}[Three‑point identity for $L$‑smooth functions]\label{lem:3pt}
For any $x,y\in\R^{d}$,
\[
\langle\nabla f(x),x-y\rangle\ge f(x)-f(y)-\frac{L}{2}\|x-y\|^{2}.
\]
\end{lemma}
\begin{proof}
Directly from the $L$‑smoothness inequality in Assumption \ref{ass:convex}. \hfill\qedsymbol
\end{proof}

\begin{lemma}[Bound on the second moment of the noise]\label{lem:noise}
If $\xi_{k}\sim\mathcal N(0,\mathbf I_{d})$ and independent of $\mathcal F_{k}$, then
\[
\E\!\big[\|\sqrt{2\alpha_{k}\tau_{k}}\;\xi_{k}\|^{2}\mid\mathcal F_{k}\big]=2\alpha_{k}\tau_{k} d .
\]
\end{lemma}
\begin{proof}
$\E\|\xi_{k}\|^{2}=d$ for a standard Gaussian in $\R^{d}$. Multiplying by the scalar $2\alpha_{k}\tau_{k}$ gives the result. \hfill\qedsymbol
\end{proof}

\section*{Main Proof (Step‑by‑Step Derivation)}
We prove Theorem \ref{thm:main}.  
Let $x^{\star}$ be any minimiser of $f$. Define the filtration $\mathcal F_{k}=\sigma(x_{0},\xi_{0},\dots,\xi_{k-1})$.  

\subsection*{Step 1: Expansion of the squared distance}
\begin{align}
\|x_{k+1}-x^{\star}\|^{2}
&=\big\|x_{k}-\alpha_{k}g_{k}+\sqrt{2\alpha_{k}\tau_{k}}\;\xi_{k}-x^{\star}\big\|^{2}\nonumber\\
&=\|x_{k}-x^{\star}\|^{2}
-2\alpha_{k}\langle g_{k},x_{k}-x^{\star}\rangle
+ \alpha_{k}^{2}\|g_{k}\|^{2}
+2\sqrt{2\alpha_{k}\tau_{k}}\;\langle\xi_{k},x_{k}-x^{\star}-\alpha_{k}g_{k}\rangle
+2\alpha_{k}\tau_{k}\|\xi_{k}\|^{2}. \tag{3}
\end{align}
[Step 1]

\subsection*{Step 2: Take conditional expectation}
Condition on $\mathcal F_{k}$. Using $\E[\xi_{k}\mid\mathcal F_{k}]=0$ and independence,
\[
\E\!\big[\langle\xi_{k},x_{k}-x^{\star}-\alpha_{k}g_{k}\rangle\mid\mathcal F_{k}\big]=0.
\]
Apply Lemma \ref{lem:noise} to the last term. Hence,
\begin{align}
\E\!\big[\|x_{k+1}-x^{\star}\|^{2}\mid\mathcal F_{k}\big]
&=\|x_{k}-x^{\star}\|^{2}
-2\alpha_{k}\langle g_{k},x_{k}-x^{\star}\rangle
+\alpha_{k}^{2}\E\!\big[\|g_{k}\|^{2}\mid\mathcal F_{k}\big]
+2\alpha_{k}\tau_{k}d . \tag{4}
\end{align}
[Step 2]

\subsection*{Step 3: Decompose $\|g_{k}\|^{2}$}
Write $g_{k}=\nabla f(x_{k})+ \delta_{k}$ where $\delta_{k}:=g_{k}-\nabla f(x_{k})$ satisfies $\E[\delta_{k}\mid\mathcal F_{k}]=0$ and $\E[\|\delta_{k}\|^{2}\mid\mathcal F_{k}]\le\sigma^{2}$ (Assumption \ref{ass:variance}). Then
\begin{align}
\E\!\big[\|g_{k}\|^{2}\mid\mathcal F_{k}\big]
&=\|\nabla f(x_{k})\|^{2}+ \E\!\big[\|\delta_{k}\|^{2}\mid\mathcal F_{k}\big]
+2\langle\nabla f(x_{k}),\E[\delta_{k}\mid\mathcal F_{k}]\rangle \nonumber\\
&\le \|\nabla f(x_{k})\|^{2}+\sigma^{2}. \tag{5}
\end{align}
[Step 3]

\subsection*{Step 4: Use $L$‑smoothness to bound $\|\nabla f(x_{k})\|^{2}$}
From $L$‑smoothness and convexity,
\[
\|\nabla f(x_{k})\|^{2}\le 2L\big(f(x_{k})-f(x^{\star})\big). \tag{6}
\]
(Proof: $L$‑smoothness gives $f(y)\le f(x)+\langle\nabla f(x),y-x\rangle+\frac{L}{2}\|y-x\|^{2}$; set $y=x^{\star}$ and rearrange.)  
[Step 4]

\subsection*{Step 5: Apply Lemma \ref{lem:3pt}}
Lemma \ref{lem:3pt} yields
\[
\langle\nabla f(x_{k}),x_{k}-x^{\star}\rangle\ge f(x_{k})-f(x^{\star})-\frac{L}{2}\|x_{k}-x^{\star}\|^{2}. \tag{7}
\]
[Step 5]

\subsection*{Step 6: Plug (5)–(7) into (4)}
Using (5) for the third term and (7) for the inner product,
\begin{align}
\E\!\big[\|x_{k+1}-x^{\star}\|^{2}\mid\mathcal F_{k}\big]
&\le \|x_{k}-x^{\star}\|^{2}
-2\alpha_{k}\Big(f(x_{k})-f(x^{\star})-\frac{L}{2}\|x_{k}-x^{\star}\|^{2}\Big) \nonumber\\
&\quad +\alpha_{k}^{2}\big(2L\big(f(x_{k})-f(x^{\star})\big)+\sigma^{2}\big)
+2\alpha_{k}\tau_{k}d . \tag{8}
\end{align}
[Step 6]

\subsection*{Step 7: Rearrange to isolate the sub‑optimality term}
Collect the terms involving $f(x_{k})-f(x^{\star})$:
\[
\big(2\alpha_{k}-2L\alpha_{k}^{2}\big)\big(f(x_{k})-f(x^{\star})\big)
\le \|x_{k}-x^{\star}\|^{2}-\E\!\big[\|x_{k+1}-x^{\star}\|^{2}\mid\mathcal F_{k}\big]
+L\alpha_{k}\|x_{k}-x^{\star}\|^{2}
+\alpha_{k}^{2}\sigma^{2}+2\alpha_{k}\tau_{k}d . \tag{9}
\]
[Step 7]

\subsection*{Step 8: Use the step‑size bound}
Assumption \ref{ass:schedules} ensures $\alpha_{k}\le 1/L$ for all sufficiently large $k$; therefore $2\alpha_{k}-2L\alpha_{k}^{2}\ge\alpha_{k}$. Hence
\[
\alpha_{k}\big(f(x_{k})-f(x^{\star})\big)
\le \|x_{k}-x^{\star}\|^{2}-\E\!\big[\|x_{k+1}-x^{\star}\|^{2}\mid\mathcal F_{k}\big]
+L\alpha_{k}\|x_{k}-x^{\star}\|^{2}
+\alpha_{k}^{2}\sigma^{2}+2\alpha_{k}\tau_{k}d . \tag{10}
\]
[Step 8]

\subsection*{Step 9: Take full expectation and sum}
Take expectation over $\mathcal F_{k}$ on both sides of (10) and sum from $k=0$ to $K-1$:
\begin{align}
\sum_{k=0}^{K-1}\alpha_{k}\E\!\big[f(x_{k})-f(x^{\star})\big]
&\le \|x_{0}-x^{\star}\|^{2}
+L\sum_{k=0}^{K-1}\alpha_{k}\E\!\big[\|x_{k}-x^{\star}\|^{2}\big]
+\sigma^{2}\sum_{k=0}^{K-1}\alpha_{k}^{2}
+2d\sum_{k=0}^{K-1}\alpha_{k}\tau_{k} . \tag{11}
\end{align}
[Step 9]

\subsection*{Step 10: Bounding the distance term}
Since the iterates remain in a bounded region under the diminishing step sizes (standard argument from Robbins–Monro), there exists $R^{2}\ge\|x_{k}-x^{\star}\|^{2}$ for all $k$. Hence
\[
L\sum_{k=0}^{K-1}\alpha_{k}\E\!\big[\|x_{k}-x^{\star}\|^{2}\big]\le L R^{2}\sum_{k=0}^{K-1}\alpha_{k}. \tag{12}
\]
We keep the bound in the final result as $\|x_{0}-x^{\star}\|^{2}$ because $R^{2}$ can be replaced by the initial distance using the telescoping nature of (8). Indeed, dropping the non‑negative term $L\alpha_{k}\|x_{k}-x^{\star}\|^{2}$ from (8) gives a simpler inequality:
\[
\alpha_{k}\big(f(x_{k})-f(x^{\star})\big)
\le \|x_{k}-x^{\star}\|^{2}-\E\!\big[\|x_{k+1}-x^{\star}\|^{2}\big]
+\alpha_{k}^{2}\sigma^{2}+2\alpha_{k}\tau_{k}d .
\]
Summing and telescoping yields
\[
\sum_{k=0}^{K-1}\alpha_{k}\E\!\big[f(x_{k})-f(x^{\star})\big]
\le \|x_{0}-x^{\star}\|^{2}
+\sigma^{2}\sum_{k=0}^{K-1}\alpha_{k}^{2}
+2d\sum_{k=0}^{K-1}\alpha_{k}\tau_{k}. \tag{13}
\]
[Step 10]

\subsection*{Step 11: From weighted sum to average iterate}
Define $S_{K}:=\sum_{k=0}^{K-1}\alpha_{k}$. By convexity of $f$,
\[
f(\bar x_{K})\le\frac{1}{S_{K}}\sum_{k=0}^{K-1}\alpha_{k}f(x_{k}),
\qquad 
\bar x_{K}:=\frac{1}{S_{K}}\sum_{k=0}^{K-1}\alpha_{k}x_{k}.
\]
Taking expectation and using (13):
\[
S_{K}\,\E\!\big[f(\bar x_{K})-f(x^{\star})\big]
\le \|x_{0}-x^{\star}\|^{2}
+\sigma^{2}\sum_{k=0}^{K-1}\alpha_{k}^{2}
+2d\sum_{k=0}^{K-1}\alpha_{k}\tau_{k}. \tag{14}
\]
Dividing by $S_{K}$ gives exactly the bound \eqref{2}. \hfill\qedsymbol

\section*{Rate Derivation (With Constants)}
Assume the concrete schedules
\[
\alpha_{k}= \frac{c}{\sqrt{k+1}},\qquad \tau_{k}= \frac{c_{\tau}}{k+1},
\]
with $c,c_{\tau}>0$. Then
\[
S_{K}=c\sum_{k=0}^{K-1}\frac{1}{\sqrt{k+1}}
\ge 2c(\sqrt{K}-1),\qquad 
\sum_{k=0}^{K-1}\alpha_{k}^{2}=c^{2}\sum_{k=0}^{K-1}\frac{1}{k+1}\le c^{2}(1+\log K),
\]
\[
\sum_{k=0}^{K-1}\alpha_{k}\tau_{k}=c c_{\tau}\sum_{k=0}^{K-1}\frac{1}{(k+1)^{3/2}}\le 2c c_{\tau}.
\]
Plugging into \eqref{2} yields
\[
\E\!\big[f(\bar x_{K})-f(x^{\star})\big]
\le 
\frac{\|x_{0}-x^{\star}\|^{2}}{2c(\sqrt{K}-1)}
+\frac{L\sigma^{2}c(1+\log K)}{2K}
+\frac{2L d c c_{\tau}}{K}.
\]
All three terms are $\mathcal O(1/\sqrt{K})$, establishing the claimed rate.

\section*{Special Cases}
\begin{enumerate}[nosep]
\item \textbf{Exact Gradient, No Noise.} If $g_{k}=\nabla f(x_{k})$ and $\tau_{k}=0$, the bound reduces to the classical SGD result
\[
\E[f(\bar x_{K})-f(x^{\star})]\le\frac{\|x_{0}-x^{\star}\|^{2}}{2\sum_{k=0}^{K-1}\alpha_{k}}.
\]
\item \textbf{Constant Temperature.} Setting $\tau_{k}\equiv\tau>0$ gives a residual term $L d\tau$ in the limit, reflecting the stationary bias of SGLD for sampling.
\item \textbf{Strongly Convex $f$.} If $f$ is $\mu$‑strongly convex, the same derivation (with Lemma \ref{lem:3pt} replaced by the strong convexity inequality) yields an $\mathcal O(\log K/K)$ rate for the averaged iterate and an $\mathcal O(1/K)$ rate for the last iterate.
\end{enumerate}

\end{document}